\chapter{Two Protocols, Two Relocations}
\label{ch:protocols}

This chapter sets Delphi and Sapience side by side. It describes how each matches a
position and how each resolves a question, traces what information each design makes
visible, to whom, and at what moment, and then asks what the verification machinery around
each one actually guarantees. The argument the chapter builds is that the two protocols
make opposite choices at both points, that neither choice removes strategic informational
advantage, and that each choice moves it somewhere else. The chapter closes by separating
three properties that are routinely conflated in discussion of AI settlement, and by
reading Delphi's own architectural history as evidence that the separation matters.

Protocol documentation changes quickly. Every factual claim here is attributed to a primary
source with the date it was read; the access dates are collected in the repository. Where a
claim could not be established from a primary source, the text says so rather than
smoothing over it.

\section{Delphi: pari-mutuel execution, AI-judged resolution}
\label{sec:delphi}

Delphi is an open information market platform built on Gensyn's decentralised AI
infrastructure. It launched on the Gensyn testnet in December 2025 and on the Gensyn network
itself on 22 April 2026, and anyone can launch a market \cite{gensyn2026delphilaunch}.
Gensyn's SDK documentation refers to a mainnet deployment \cite{gensyn2026sdk}; the launch
post does not use that word, so no mainnet launch date is asserted here.

\subsection{Pricing}

Gensyn rejected both dominant market designs explicitly, and their stated reasoning is
worth recording because it explains the design that resulted. Gensyn's position is that a
logarithmic market scoring rule \cite{hanson2007lmsr} requires initial liquidity and that
the party providing it often loses money, which they judge tolerable when a platform
absorbs that loss across a curated set of markets but unworkable under permissionless
creation; and that an order book requires a counterparty on the other side of every trade,
which in a thin market is often nobody \cite{gensyn2026delphi}. An automated market maker
of the constant-function kind is a third possibility that Gensyn do not discuss, and which
Sapience also declines in favour of the auction described in
Section~\ref{sec:sapience}.

Delphi therefore prices with a dynamic pari-mutuel cost function
\cite{pennock2004dpm,gensyn2026delphi}. In a classical pari-mutuel market all stakes enter
a pool and a single price is fixed at close, so a trader early in a market's life is
guessing at the price they will eventually receive. Pennock's dynamic variant keeps the
pari-mutuel property that the institution bears no risk, since payouts come from the pool,
and adds continuous pricing during trading, so that a trader can react to new information
and can sell before the event resolves; the price derives from the money wagered on each
outcome and the shares outstanding against it \cite{pennock2004dpm}. As demand accumulates on one side, the marginal cost of buying that
side rises and the implied payout adjusts. The market creator chooses any non-zero initial
liquidity and retains a position in their own market \cite{gensyn2026delphi}.

Two deployments must be distinguished. Gensyn's SDK documentation states that Delphi application markets use dynamic
pari-mutuel pricing, and that automated settlement changes settlement and routing without
changing Delphi markets to a logarithmic scoring rule. The separate Gensyn agent trading
competition, on \texttt{competition-testnet}, does use LMSR contracts through the same SDK
trading interfaces \cite{gensyn2026sdk}.

Each market is deployed through a factory as its own non-upgradeable contract. Gensyn can
moderate what appears in its application but cannot rewrite a market after launch. The
contracts do not distinguish humans from agents: anyone with a wallet may buy, sell and
redeem against the same liquidity \cite{gensyn2026delphi}.

\subsection{Settlement}

Gensyn's stated position is that human resolution committees are slow, capturable and
expensive; that optimistic-oracle staking designs are slow by construction and
capital-inefficient at the long tail; and that AI judges are fast and cheap but, in a
closed system, ask the user to trust the operator. Their conclusion is that the AI judge is
the right primitive and that the work lies in making it trustworthy, with a model reading
real-world sources and judging what has already happened rather than forecasting what is to
come \cite{gensyn2026delphi}.

Settlement is offered in tiers, so a creator may trade speed against verifiability.
Foundation models settle quickly and broadly, chosen from a curated list. Open-source models
running inside Gensyn's Reproducible Execution Environment produce a receipt that can be
independently rerun to verify the answer, settling more slowly in exchange for that
property. Hosted variants let creators outsource the reproducible tier to dedicated
infrastructure providers \cite{gensyn2026delphi}.

The detail that matters for this project is who writes the settlement procedure. In the
design documented in May 2026, creators should write the prompts that settle their own
markets, on the reasoning that a creator understands their question and its edge cases
better than the platform does. Gensyn report
that testnet experience taught them that the cheapest way to obtain a wrong settlement is a
vague prompt, and they responded with category-aware templates across sports, politics,
crypto, economics, science and culture, one template per settlement path, that supply
structure while the creator supplies the substance. Market questions also pass a keyword
blocklist and a semantic check \cite{gensyn2026delphi}.

Gensyn are candid about the limit of this. Their own framing is that none of it makes the
AI judge infallible; it makes the failure modes legible, so that a bad settlement traces
back to a specific prompt, model and run rather than to an opaque operator decision, and on
the reproducible tier anyone may rerun the call \cite{gensyn2026delphi}. Legibility and
reproducibility are real properties. Section~\ref{sec:threelayer} argues that neither is
the property a disputed settlement actually needs.

The current deployment differs from the design just described. Gensyn's SDK documentation
states that testnet and mainnet use automated-settlement gateways by default, that some
legacy markets remain on the earlier deployment, and that the client resolves a market's
owning gateway before each market-scoped operation. Automated settlement changed gateway
addresses and market routing, and market statuses now include a failed state that requires
liquidation rather than redemption \cite{gensyn2026sdk}. Section~\ref{sec:evolution} treats
this migration as evidence rather than as background.

\section{Sapience: RFQ execution, optimistic-oracle resolution}
\label{sec:sapience}

Sapience is a prediction market protocol, described by its own repository as MIT-licensed
and fully open source, in which users stake collateral on
the outcomes of future events. Positions are matched by a request-for-quote auction and
settled on-chain by an escrow contract, with resolution delegated to an optimistic oracle
\cite{sapience2026relayer}.

\subsection{Matching}

Sapience uses neither an order book nor an automated market maker. The auction relayer is a
WebSocket service brokering a request-for-quote flow between two roles. A predictor
initiates an auction by broadcasting their picks and the collateral they wish to commit,
authenticated by an EIP-712 \texttt{AuctionIntent} signature. A counterparty, typically a
vault bot, watches auctions and submits a bid carrying the EIP-712 \texttt{MintApproval}
signature that allows the predictor to mint the on-chain position
\cite{sapience2026relayer}.

Figure~\ref{fig:rfqflow} shows the sequence. Bids carry short expiries and the auction
resolves in seconds \cite{sapience2026docs}.

\begin{figure}[htbp]
\centering
\begin{tikzpicture}[
  font=\small,
  lifeline/.style={draw=black!35, dashed, thick},
  actor/.style={draw, rounded corners=2pt, fill=black!4, inner sep=5pt, minimum width=27mm},
  msg/.style={-{Latex[length=2mm]}, thick},
]
\node[actor] (P) at (0,0) {Predictor};
\node[actor] (R) at (5.4,0) {Relayer};
\node[actor] (C) at (11.0,0) {\begin{tabular}{c}Every connected\\client\end{tabular}};

\foreach \x in {0,5.4,11.0} \draw[lifeline] (\x,-0.55) -- (\x,-7.4);

\draw[msg] (0,-1.2) -- node[above,font=\scriptsize] {\texttt{auction.start}} (5.4,-1.2);
\draw[msg] (5.4,-2.0) -- node[above,font=\scriptsize] {\texttt{auction.ack}} (0,-2.0);

% the disclosure moment, banded so it reads as one step
\begin{scope}[on background layer]
  \fill[black!7, rounded corners=3pt] (4.6,-4.35) rectangle (11.9,-2.7);
\end{scope}
\draw[msg, very thick] (5.4,-3.1) -- node[above,font=\scriptsize\bfseries] {\texttt{auction.started}} (11.0,-3.1);
\node[font=\scriptsize\itshape, text=black!65, text width=58mm, align=center, anchor=north, fill=black!7, inner sep=1.5pt]
  at (8.2,-3.3) {every pick with its predicted outcome,\\the position size, the predictor's address};
\node[font=\scriptsize\itshape, text=black!65, text width=58mm, align=center, anchor=north, fill=black!7, inner sep=1.5pt]
  at (8.2,-3.95) {no allowlist, no subscription, no filtering};

\draw[msg] (11.0,-5.0) -- node[above,font=\scriptsize] {\texttt{bid.submit}} (5.4,-5.0);
\draw[msg] (5.4,-5.8) -- node[above,font=\scriptsize] {\texttt{auction.bids}} (0,-5.8);
\draw[msg] (0,-6.9) -- node[above,font=\scriptsize] {\texttt{mint()} on-chain, both signatures} (5.4,-6.9);
\end{tikzpicture}
\caption{The Sapience request-for-quote flow. The shaded step is where a forecaster's belief
becomes visible: it precedes any fill, it reaches every connected client rather than a
selected set of market makers, and it carries the complete request. Everything before it is
private to the predictor and everything after it is a response to what it revealed.}
\label{fig:rfqflow}
\end{figure}

The relayer never takes custody. Sapience's stated trust model is that the relayer only
routes messages and can neither forge bids nor modify position sizes nor steal funds; that
responder signatures are verified on-chain, so a bid is worthless unless correctly signed;
that the immutable audited contract is the sole authority over collateral and settlement;
and that liquidity vaults hold pooled capital whose manager decides which auctions to bid on
but cannot withdraw funds \cite{sapience2026docs}. Section~\ref{sec:execflow} returns to
what this trust model does and does not cover.

\subsection{Combinations}

A prediction is one or more picks together with a position size, where a pick is a question
paired with an expected resolution. Multi-pick predictions combine several into a single
position that pays only if every pick resolves correctly. For independent picks the implied
probability is the product of the individual probabilities, so a combination
carries a lower combined probability and a correspondingly higher payout
\cite{sapience2026docs}. Combinations therefore express compound forecasts that a
single-event venue cannot. They also change what a broadcast reveals, and in the strongest
way available to this argument. A single-pick request discloses a marginal belief about one
event. A multi-pick request discloses the direction the forecaster takes on every leg at
once, which is a statement about the joint structure of their beliefs rather than a set of
independent statements: it identifies which combination of outcomes they will stake on
together. Section~\ref{sec:execflow} treats the broadcast payload accordingly.

\subsection{Resolution}

Sapience settles through UMA's optimistic oracle. After a market's end time the deployer
submits a resolution with a bond. During a challenge period anyone may dispute by posting a
bond of their own. If the assertion is undisputed, participants withdraw according to it. If
it is disputed, the question passes to a vote of UMA stakers, and the losing side forfeits
its bond \cite{sapience2026docs,uma2026oracle}.

UMA's documentation sets out the voting mechanism. Stakers vote in a twenty-four hour commit
phase followed by a twenty-four hour reveal phase. A dispute resolves when a minimum
sixty-five per cent majority of staked UMA is cast for a single outcome. Slashing and
rewards are described as incentivising stakers to vote for the most correct answer, and
secret ballots are described as preventing lazy voters from copying others and dishonest
voters from coordinating on an incorrect result \cite{uma2026oracle}.

\section{Information flow at execution}
\label{sec:execflow}

The two designs expose a trader's intent at opposite moments. This section establishes what
is exposed, to whom, and when. Table~\ref{tab:execflow} summarises it.

\begin{table}[htbp]
\centering
\begin{tabularx}{\textwidth}{p{0.22\textwidth}p{0.34\textwidth}X}
\toprule
\textbf{Question} & \textbf{Delphi} & \textbf{Sapience} \\
\midrule
When is intent visible & After the trade, through its effect on the quoted price & Before any fill, in the request itself \\
What is visible & Direction and size, aggregated into a price movement across the share supply & The complete pick set, with the predicted outcome on every leg, and the position size \\
Who can see it & Any on-chain observer, after the fact & Every client connected to the relayer, at broadcast time \\
Is the trader linkable across requests & By the transacting address, on-chain & By the predictor address, carried in the broadcast payload \\
Can a viewer act before the trader is filled & Only through ordinary transaction-ordering means & Yes; that is the mechanism's purpose \\
\bottomrule
\end{tabularx}
\caption{Intent exposure at execution. The row that carries the argument is the second: the
two venues do not merely differ in timing, they differ in how much of a belief is
recoverable from what they expose.}
\label{tab:execflow}
\end{table}

On Delphi, a trade is an on-chain transaction against a pari-mutuel cost function. It
changes the share supply and therefore the quoted price, so other participants learn from
the price movement after the position is established. What they observe is aggregated: the
price reflects the net of everything traded, and a single trader's belief is recoverable
only to the extent that their trade moved the price. This is not an absence of a strategic
surface. On-chain execution carries the ordinary transaction-ordering surface that the MEV
literature describes, in which a party able to sequence or observe pending transactions can
act on them before they are included \cite{daian2019flashboys,ethereummev}. The point is
narrower: Delphi does not add a pre-trade disclosure channel of its own on top of that.

On Sapience the disclosure is deliberate, complete, and prior to any fill. Three properties
of the broadcast, established from the protocol's own documentation and implementation, are
load-bearing for the rest of this report.

First, the broadcast is unfiltered. The relayer's handler context declares a function
returning all connected clients, annotated in the source as being used for the global
auction broadcast, and the broadcast loop iterates every open connection with no
subscription requirement, no responder allowlist and no filtering
\cite{sapience2026relayer}. The endpoint is public and documented in the protocol's builder guide
\cite{sapience2026builderguide}, which states that the relayer broadcasts the auction to
every connected client and describes the counterparty role as one that watches every
auction.

Second, the payload is rich. The broadcast structure carries the full array of picks, each
of which names a condition and the predicted outcome for it; the predictor's collateral,
which is the position size; and the predictor's address \cite{sapience2026relayer}. A
responder therefore observes the direction of every leg of a compound forecast, the stake
behind it, and an address, so that repeated use of the same address permits longitudinal
linking of one participant's requests.

Third, the mechanism offers no way to disclose less and still trade. The intent signature is
optional at the protocol level, but the documentation states that vault counterparties
respond with actionable bids only to signed requests, and that unsigned requests should be
treated as price discovery only \cite{sapience2026builderguide}. A forecaster who wants a
fill on a given request must disclose that request in full. This is a claim about a single
request and not about a trading strategy: a forecaster may still split a position across
several smaller requests, which discloses less per request at the cost of executing in
pieces. What is unavailable is a priced option to reveal a coarser version of the request
one is actually making.

\subsection{What follows, and what does not}

It would be easy, and wrong, to describe this as a privileged information channel. Access to
the stream is not permissioned. Any client that opens a WebSocket connection receives every
auction, so the advantage available from watching order flow attaches to the capital and
infrastructure needed to quote on it rather than to occupying a protected role. This
observation cuts against the simplest version of the argument this report makes, and
Chapter~\ref{ch:evaluation} treats it as a finding rather than as an inconvenience.

What the design does supply cheaply is linkability. Because an address appears in every
broadcast, any observer can accumulate a per-address history at no cost and from no
privileged position. One qualification is owed immediately: nothing compels a forecaster to
reuse an address, so continuity is available to an observer rather than imposed on a
discloser. Rotation is not free, since an address that has no history also has no standing
with the vaults quoting on it, but the cost is a matter of reputation rather than of
protocol, and the accumulation argument should be read as conditional on a forecaster who
does not rotate. Whether such accumulation converts into an advantage that competitive
quoting does not compensate is a separate question, which this report argues rather than
demonstrates; the simulator of Chapter~\ref{ch:method} models exposure and not cross-event
learning, and Chapter~\ref{ch:evaluation} is explicit about that boundary.

Finally, none of this is concealed by the protocol. Sapience's trust model addresses custody, forgery and
contractual authority. It does not consider what responders learn from observing requests. That is a category the trust model
does not model, rather than a flaw it hides, and this report treats it as an unmodelled
surface. The distinction matters for how the finding should be read and it matters ethically, a point set out in Section~\ref{sec:ethics-intro}.

\section{Information flow at resolution}
\label{sec:resflow}

Both protocols must convert a real-world state of affairs into a payout, and both delegate
that conversion. They delegate it to different parties, with different failure
characteristics.

On Delphi the specification originates with the market creator, who also holds a position in
the market. In the design Gensyn documented in May 2026, the creator defined the outcomes
and the data sources, wrote the settlement prompt within a category template, and selected
the settling model from a curated list \cite{gensyn2026delphi}. Whether prompt authorship
and model selection remain with the creator under the automated-settlement deployment could
not be established: Gensyn's current documentation describes no settlement mechanism, no
oracle and no prompt authorship, and the SDK pages consulted for this report are silent on
all three \cite{gensyn2026sdk}. The analysis below therefore concerns specification power
as a structural feature of creator-defined markets, and does not assert who exercises it
today. Each of the outcomes, the data sources, the prompt and the model is a degree of
freedom that shapes the resolved outcome without any of them constituting misbehaviour. A settlement question written in good
faith still embeds a reading of the event, and this report reads Gensyn's own observation, that vague
prompts are the cheapest route to a wrong settlement, as consistent with the specification
being where the outcome is decided.

On Sapience the specification is fixed when the market is created, and the contested step is
the assertion. The deployer asserts an outcome under bond; a challenger may contest it under
an equal bond; a disputed question is decided by a token-weighted vote
\cite{sapience2026relayer,uma2026oracle}. The mechanism prices the act of asserting
falsely, which is a real protection against cheap misassertion.

The two designs are also each other's stated alternative, which is worth noting because the
critique is first-party rather than this report's. Gensyn's reasoning for rejecting an
optimistic oracle, recorded in Section~\ref{sec:delphi}, is that such designs are slow by
construction and capital-inefficient at the long tail \cite{gensyn2026delphi}. Sapience
settles through exactly that mechanism. Neither party is disinterested, and the point is not
that one is right, but that each protocol's designers regard the other's resolution choice
as carrying a cost they were unwilling to pay.

The two failure modes are close to mirror images. A Delphi settlement that turns on an
ambiguous specification fails silently, quickly and cheaply: a model reads the prompt,
returns an answer, and the market pays out, with the ambiguity visible only to a participant
who reads the prompt and disagrees with it. A Sapience settlement that turns on an ambiguous
specification fails loudly, slowly and expensively: someone must post a bond, wait out the
liveness period, and win a vote. Neither mechanism removes the ambiguity. One absorbs it
without friction and the other converts it into a costly contest.

\subsection{An argument about optimistic-oracle voting}
\label{sec:umaargument}

The following is this report's own argument about a documented mechanism, not a reported
result, and it is labelled as such because no independent scholarly treatment of the point
was located. The claim that optimistic-oracle voters are incentivised to vote with the
expected majority rather than with the truth is widely repeated and, so far as could be
established, rarely sourced. Rather than cite an authority that does not exist, the argument
is made here from UMA's own description of the mechanism.

A vote that pays for agreement with the resolved majority is a coordination device. Such a
device tracks truth for exactly as long as the truth is the obvious thing for other voters
to choose. On a question whose wording is unambiguous, it is; the correct answer is the
natural focal point and the mechanism works as intended. On a question whose wording is
genuinely contested, the focal point is whatever voters expect other voters to select, and
there is no guarantee that this coincides with the best reading of the specification. UMA's
defences are well designed against the failure they target: secret ballots frustrate
copying and coordination, a sixty-five per cent threshold raises the share of stake an
attacker must control, and the delay on unstaking exposes an attacker to the depreciation
their own attack would cause \cite{uma2026oracle}. Each of these makes it harder for voters
to agree on a falsehood. None of them makes it easier for voters to agree on a meaning. The
residual is not dishonesty and it is not capture; it is ambiguity, and the mechanism has no
instrument that addresses it.

\section{Threat models}
\label{sec:threat}

The two preceding sections traced what each protocol exposes. This one asks who is in a
position to act on it, and what the protocol does about that. Table~\ref{tab:threat} lists
only the actors whose powers are not already covered by Sections~\ref{sec:execflow}
and~\ref{sec:resflow}; traders, ordinary observers and dispute voters are omitted because
their positions were established there. The table is descriptive, and naming a power is not
an allegation that anyone exercises it.

\begin{table}[htbp]
\centering
\small
\begin{tabularx}{\textwidth}{p{0.15\textwidth}p{0.21\textwidth}p{0.27\textwidth}X}
\toprule
\textbf{Actor} & \textbf{Power held} & \textbf{Constrained by} & \textbf{Not constrained} \\
\midrule
\multicolumn{4}{l}{\emph{Delphi}} \\
Market creator & Defines outcomes, data sources, settlement prompt and model; holds a position & Category templates; abuse filter; non-upgradeable market contract & Whether the specification is a neutral rendering of the question \\
Settlement operator & Executes the settlement call & Reproducible receipts on that tier; under automated settlement, gateways the creator does not operate & Model choice, timing, and any external source the judge reads \\
\midrule
\multicolumn{4}{l}{\emph{Sapience}} \\
Relayer & Routes messages & Cannot forge bids, alter sizes or take custody; signatures verified on-chain & Observes complete order flow \\
Responder or vault & Quotes on observed requests & Signature verification; the predictor's freedom to reject a bid & What it learns from the request, and from linking requests by predictor address \\
Market deployer & Asserts the resolution under bond & Bond, challenge window, dispute vote & Whether the specification was neutral \\
\bottomrule
\end{tabularx}
\caption{Actors whose powers are not fully covered by the information-flow analysis. The
final column is where the argument lives.}
\label{tab:threat}
\end{table}

Read down the third column and the protocols look well defended, because they are. Every
constraint listed is real, most are enforced on-chain or cryptographically, and between them
they close the failures a designer would name first: a relayer that forges bids, an operator
that substitutes a settlement result, a deployer that asserts a false outcome without cost,
a creator who rewrites a market after people have traded it.

Read down the fourth column and a pattern appears that neither protocol's documentation
draws attention to. The unconstrained powers are not failures of enforcement. They are
things the enforcement was never aimed at, and they divide into two kinds.

The first kind is observation. The relayer cannot steal, forge or alter, and it sees every
request; a responder is bound by signature verification and by the predictor's freedom to
walk away, and it learns whatever the request contains. Sapience's trust model is accurate
about custody, forgery and authority, and simply does not treat observation as a harm to be
modelled. That is a category the model omits rather than a flaw it conceals, which is why
Section~\ref{sec:execflow} describes it as an unmodelled surface.

The second kind is specification. A Delphi market creator holds a position and writes the
criteria that decide it; a Sapience deployer asserts an outcome under a specification fixed
earlier. The constraints in both cases operate on execution and on assertion, and neither
touches the content of the specification itself.

What the two kinds have in common is that no participant exercising them is misbehaving. A
responder pricing a request it was sent is doing the job the mechanism created for it. A
creator writing a settlement prompt is doing what the platform asks. The protections
enumerated in the third column are aimed at dishonesty, and both residuals survive them
because neither residual is dishonesty. Section~\ref{sec:threelayer} takes up the second of
them, which is the more tractable to state precisely.

\section{Three verification properties, and what each bounds}
\label{sec:threelayer}

Verifiability is not one property. This section separates the properties that the mechanisms
described above actually deliver, because they answer different failure modes and are not
interchangeable. The separation is this report's analytical device rather than a reported
finding, and it is the chapter's principal contribution.

\subsection*{Reproducibility: did this model, on these inputs, produce this output?}

Reproducibility is delivered by rerunning an open-source model inside a controlled
environment and comparing the result. It bounds divergence between runs, and it converts a
settlement from an assertion into a computation a third party can repeat.

Gensyn are precise about where it stops. A receipt proves that a given model produced a
given output from a given prompt, and does not verify that the information in the prompt is
accurate; and receipts do not execute tools, so they do not prove that an external tool
returned accurate, complete or unchanged data \cite{gensyn2026receipts}. The second
limitation is not marginal for a settlement judge whose entire task is reading real-world
sources: it means the guarantee covers the reasoning and not the evidence.

\subsection*{Execution attestation: did this call actually happen as described?}

Attestation is delivered by executing a model call inside a trusted execution environment
and attesting the transcript on a system outside the application's control. It
bounds prompt substitution, selection among responses, and retrying until a preferred answer
returns.

It is the necessary instrument for closed frontier models, and the reason is technical
rather than philosophical. Such models are non-deterministic: calling one a second time with
the same prompt generally yields a different answer, so an inference cannot be proved after
the fact by repeating the call \cite{truebit2026orchestration}. Reproducibility is therefore
available for open models under a controlled environment, and attestation is the
corresponding answer for closed ones, whose outputs can be witnessed but not re-derived. A
stack offering one is not thereby offering the other, and this is the point at which
treating verifiability as a single property produces a wrong conclusion about what a
deployment guarantees.

\subsection*{Semantic validity: was the question a fair rendering of the event?}

Semantic validity asks whether the question, the evidence and the resolution criterion were
a fair and unambiguous rendering of what happened. Both protocols do address it, and it
would be wrong to say otherwise: Delphi supplies category-aware templates, definitions and
edge-case prompts, and screens questions; Sapience fixes the specification before trading
and lets a bonded challenger contest an assertion made under it. These are real instruments
and they improve specifications.

\textbf{What no mechanism in either protocol does is guarantee semantic validity, and
verification in particular cannot.} That is not an omission in either protocol's
engineering; it is a statement about what verification is. Reproducibility and attestation
are both claims about a computation, and semantic validity is a claim about meaning.
Establishing that a procedure ran exactly as committed says nothing about whether the
committed procedure was the right one, and no amount of further verification changes that,
because the residual is not of the same kind as the thing being verified. Templates and
dispute processes reduce the residual; they do not close it, and neither is a verification
mechanism in the sense the other two rows of Table~\ref{tab:threelayer} use the word.

Table~\ref{tab:threelayer} summarises the three.

\begin{table}[htbp]
\centering
\begin{tabularx}{\textwidth}{p{0.26\textwidth}p{0.28\textwidth}X}
\toprule
\textbf{Property} & \textbf{Mechanism} & \textbf{What it leaves open} \\
\midrule
Reproducibility & Rerun an open model in a controlled environment & Whether the prompt was accurate, and whether external sources were \\
Execution attestation & TEE-executed call, transcript attested off the operator's infrastructure & Whether the committed specification was neutral \\
Semantic validity & Specification templates, precommitted criteria, and bonded contestation. \textbf{No verification mechanism} & Whether a reading is the uniquely correct one; the residual is reduced, not closed \\
\bottomrule
\end{tabularx}
\caption{The three properties in summary; the argument is in the text above. The third row is
the chapter's central observation: semantic validity is addressed by procedure and never by
verification, so it is the one property the other two mechanisms cannot deliver.}
\label{tab:threelayer}
\end{table}

\subsection*{Attested execution in practice}

Truebit describe an unnamed permissionless AI-settled prediction-market deployment in which
a dynamic oracle detects a market entering settlement and triggers the workflow without a
manual trigger, settlement calls to frontier models run inside a trusted execution
environment with the exact prompt and response captured in an attested transcript on
a system outside the application's control, and deterministic logic then runs
across multiple verifier nodes, with each step producing a receipt
\cite{truebit2026orchestration}. The threat they name as the motivation is precisely the one
this chapter has been tracing: whether the party with the most at stake in an outcome could
have shaped it, by adjusting a prompt, selecting among responses, or retrying until a
preferred answer returned \cite{truebit2026orchestration}.

Two things must be said carefully. The first is identification. The Truebit case study does
not name Delphi or Gensyn, and Gensyn's documentation does not name Truebit: the settlement
methods page names no oracle relayer and describes no settlement mechanism at all
\cite{gensyn2026sdk,truebit2026orchestration}. Truebit's description is consistent with an
automated settlement path of the kind Gensyn's SDK routes to, but since Gensyn document no
settlement mechanism there is nothing on their side to compare it against, so no
correspondence can be established and this report does not assert one.

The second is what attestation achieves. Attestation mitigates execution-time manipulation
of a committed resolution procedure. It does not establish that the underlying market
specification is semantically neutral. On Truebit's own description of the workload, the
creator still defines the market's outcomes and data sources, which is specification power
exercised before any execution occurs and therefore untouched by any guarantee about how
that execution proceeded \cite{truebit2026orchestration}.

\section{Delphi's evolution as evidence for the framework}
\label{sec:evolution}

The separation set out above is not only analytical. Delphi has moved along it, and Gensyn's
own documentation records the move. Under the arrangement documented in May 2026 the
settling party held both semantic and execution discretion. Under automated-settlement
gateways the creator no longer executes the settlement, and Gensyn publish nothing to
indicate that specification discretion moved with it. What first-party sources establish is
the removal, not what remains: the current settlement procedure, prompt authorship and model
selection are not described anywhere in Gensyn's documentation, as
Section~\ref{sec:delphi} records.


The three settlement tiers of Section~\ref{sec:delphi} are concurrent options rather than
successive designs: Gensyn describe them as a path that lets the creator trade speed for
verifiability \cite{gensyn2026delphi}. Reproducible execution was therefore available
\emph{within} the creator-defined arrangement rather than replacing it, which is why it
appears as a property in this chapter and not as a period.

What the sources do establish is the transition between the two rows. Gensyn's SDK
documentation states that testnet and mainnet use automated-settlement gateways by default,
that some legacy markets remain on the earlier deployment, and that automated settlement
changed gateway addresses and market routing \cite{gensyn2026sdk}. Gensyn do not quantify
what share of markets remains on the earlier deployment, and this report does not estimate
it. The direction of travel is from settlement the creator executed to settlement the
creator does not execute.

A third architecture is documented in the same period, and it is worth setting beside these
two without being folded into them. Truebit describe a deployment in which settlement calls
to frontier models run inside a trusted execution environment with the prompt and response
captured in an attested transcript on a system outside the application's control, and
contrast it with markets that rely on a whitelisted set of trusted creators
who run the settlement model themselves, which they characterise as a reasonable way to
bootstrap a platform and not a way to scale one \cite{truebit2026orchestration}. That
client is unnamed, and Gensyn's documentation describes no settlement mechanism against
which the description could be matched, so the attested architecture is treated here as a
documented design point of the same period rather than as Delphi's own third state.

The reading this report takes is that the change removes a genuine class of privilege, and
that it is a serious engineering response to a real threat rather than a concession that
something was wrong.

What makes the sequence analytically useful is what survives it. In both documented states
the settling arrangement leaves semantic discretion somewhere, because somebody must decide
what the question means, which sources count as evidence, and what will be treated as the
criterion for resolution. Removing the creator's execution discretion does not touch that,
and neither would attestation, because verification is a statement about a computation and
the residual is a statement about a specification.

The claim that the evolution is \emph{evidence} for the separation, rather than merely an
instance of it, rests on one point and it is worth stating rather than assuming. The
framework was not imposed on the sequence: a designer with every incentive to close the
whole problem closed execution discretion and published nothing to indicate that
specification discretion was closed with it. That is an argument from silence, and the
silence has a competing explanation, since Gensyn document no settlement mechanism at all;
but the choice that is documented was made independently of this report's categories, and
it is what gives the sequence purchase beyond restating them.

This is the sense in which the chapter's title should be read. Delphi and Sapience relocate
strategic informational advantage rather than eliminating it. At execution, Delphi's
pari-mutuel pool relocates it to post-trade inference and the ordinary transaction-ordering
surface, while Sapience's request-for-quote auction relocates it to a pre-trade broadcast
that is complete, address-linked and unpriced. At resolution, Delphi relocates it into a
specification interpreted by a machine under strengthening execution guarantees, while
Sapience relocates it into a bonded assertion contested by a coordination mechanism. In all
four cases the mechanism that works is aimed at dishonesty, and what remains after it has
worked is meaning.

\section{What this chapter established}

This chapter established the following, and it is worth separating what was read from what
was reasoned. ``Demonstrated'' is reserved throughout this report for what the experiment
measures; what a document states is \emph{established}, which is a different act.

Established from primary sources:

\begin{itemize}
  \item Delphi prices with a dynamic pari-mutuel cost function in production, while a
    separate agent-competition deployment uses a logarithmic scoring rule.
  \item In the design Gensyn documented in May 2026, market creators defined the outcomes,
    data sources, settlement prompt and settling model while holding a position.
  \item Delphi has since migrated to automated-settlement gateways, with some legacy markets
    remaining on the earlier deployment.
  \item The Sapience relayer broadcasts each request to every connected client, carrying the
    complete pick set with the predicted outcome on each leg, the position size and the
    predictor's address, and offering no partial-disclosure path to anyone intending to trade.
\end{itemize}

Argued, not demonstrated:

\begin{itemize}
  \item That access to the Sapience broadcast is unpermissioned, so the resulting advantage
    attaches to capital and infrastructure rather than to a protected role.
  \item That an address in every broadcast makes accumulation cheap. This does not
    establish that accumulation converts into uncompensated advantage.
  \item That optimistic-oracle voting is a coordination device whose defences target
    dishonesty and not ambiguity.
  \item That semantic validity, reproducibility and execution attestation are three
    properties, of which neither protocol's verification machinery guarantees the first.
\end{itemize}

Not established, and stated as such: whether prompt authorship and model selection remain
with the creator under the automated-settlement deployment, on which Gensyn's current
documentation is silent. Neither gap is bridged by inference here. Gensyn also do not
quantify how many markets remain on the legacy deployment, and no estimate is offered.

Chapter~\ref{ch:method} takes the pre-trade disclosure identified in
Section~\ref{sec:execflow} and asks what it is worth, by measuring the surplus that moves
when responders may condition on an observed request against a benchmark in which they may
not.
